% SPDX-License-Identifier: GPL-3.0-or-later
\documentclass[12pt]{article}
\usepackage[margin=1in]{geometry}
\usepackage{amsmath,amssymb,amsthm,siunitx,derivative,wrapfig,longtable,tcolorbox,listings,hyperref}
\usepackage{cleveref}
\newtheorem{example}{Example}[section]
\lstset{basicstyle=\ttfamily\small,keywordstyle=\color{blue},commentstyle=\color{green!40!black},frame=single,breaklines=true}
\DeclareSIUnit{\widget}{wd}
\newcommand{\greeting}[2][Hello]{#1, #2!}
\newcommand{\qtycm}[1]{\qty{#1}{\centi\metre}}
\begin{document}
\section{Links}\label{sec:links}
Visit \href{https://www.latex-project.org/}{the \LaTeX{} project} or
\url{https://ctan.org/pkg/siunitx}.
\hypertarget{start}{This is a return point.}
Use \hyperlink{start}{this link} or \hyperref[sec:links]{the Links section}.

\section{Numbers and units}
A distance is \qty{1.23e4}{\kilo\metre}.
An acceleration is \qty{9.81}{\metre\per\second\squared}.
Compare \qty[per-mode=fraction]{3.5}{\metre\per\second} with
\qty[locale=DE]{12345.67}{\metre}.
Custom units work too: \qty{4}{\widget}.
Small time intervals use \qty{450}{\micro\second}.
An argument-preserving wrapper gives \qtycm{10}, also inside mathematics:
$d=\qtycm{25}$.

A range is \qtyrange{1}{5}{\second}; a list is \numlist{1;2;3};
an angle is \ang{30;15;20}.
Legacy commands include \SI{10}{\kilogram} and \si{\newton\metre}.

\section{Derivatives}
\begin{align*}
  \odv{f}{x} &= 2x, & \odv[order=2]{f}{x} &= 2, \\
  \pdv{g}{t} &= \cos t, & \pdv[order={2,1}]{h}{x,y} &= 0, \\
  \odv[order=n]{f}{x} &= 0, & \mdv*{F}{t} &= 0.
\end{align*}
Differentials include $\odif{x}$, $\pdif{y}$ and $\fdif{S}$.
Slash notation is $\odv{f}/{x}$.

\section{Boxes, rules and spacing}
Fixed-width text can be aligned using \verb|\makebox[3cm][r]{Right}|:
\makebox[3cm][r]{Right}. A framed box is \fbox{Text $x^2$};
\framebox[3cm][l]{Left} keeps its chosen width.
A rule is \rule[2pt]{2cm}{0.6pt}, and \raisebox{2pt}{raised text}
shares the paragraph with \mbox{unbroken text}.
\noindent Explicit spacing: left\hspace{1em}right.

\section{Optional defaults}
\greeting{reader}
Change the first argument of \greeting[Welcome]{reader} to its default,
Hello: PasTeX removes the redundant brackets from the source.
\section{Colored boxes and theorem references}
\begin{tcolorbox}[title={An editable title with $x^2$},colback=blue!5,colframe=blue!60!black]
Both this title and the body can be edited in place. Text can be
\textbf{bold}, and mathematics such as $\frac{1}{1+x^2}$ retains its own fields.
\end{tcolorbox}
Click in the box to open its Color box tab and configure its title, colors,
corner radius or border. Empty properties inherit the preamble settings.

\begin{example}[A positive square]\label{ex:square}
If $x\in\mathbb{R}$, then $x^2\geq0$. A matrix may be written as
$A=\begin{pmatrix}1&x\\0&1\end{pmatrix}$.
\end{example}
\begin{equation}\label{eq:package-one} a^2+b^2=c^2.\end{equation}
\begin{equation}\label{eq:package-two} e^{i\pi}+1=0.\end{equation}
References include \cref{ex:square}, \cref{eq:package-one,eq:package-two},
\Crefrange{eq:package-one}{eq:package-two}, and \nameref{sec:links}.

\section{Wrapped figures}
\begin{wrapfigure}{r}{.35\textwidth}
\centering
\begin{tabular}{|c|c|}\hline$x$&$x^2$\\\hline1&1\\2&4\\\hline\end{tabular}
\caption{A small wrapped illustration.}\label{fig:wrapped}
\end{wrapfigure}
This paragraph flows beside the illustration before returning to the full
line width. Edit the caption or any entry in its small table. The width and
placement are ordinary arguments of the wrapped environment. Text entered
beside the figure is still part of this paragraph. Continue typing here to
see the paragraph reflow while the figure stays beside the text. A longer
paragraph eventually continues underneath the illustration, leaving enough
room for both its content and caption. See also \cref{fig:wrapped}.

\section{Listings}
Literal characters inside a listing remain editable as source code. Inline code such as \lstinline|WriteLn('\{Hello\}')|
is editable too. Edit a word, insert a line, or copy part of the code.
\begin{lstlisting}[language=Pascal,numbers=left]
program Greeting;
begin
  { Literal backslashes and braces stay visible. }
  WriteLn('\{Hello\}');
end.
\end{lstlisting}
The Listing tab edits its language, highlighting, font, numbering and frame.
For example, choose \verb|basicstyle=\ttfamily| for a monospace listing.

\section{Long tables}
The continuous instant preview shows the initial header and final footer
once. The TeX engine repeats continuation headers when the table spans pages.
\begin{longtable}{|l|r|}
\caption{Editable longtable cells and caption.}\label{tab:packages}\\
\hline Name&Value\\\hline
\endfirsthead
\hline Name (continued)&Value\\\hline
\endhead
\hline\multicolumn{2}{r}{Continued on the next page}\\
\endfoot
\hline
\endlastfoot
Alpha&1\\
Beta&2\\
Gamma&3\\
\hline
\end{longtable}
See \cref{tab:packages} for a table reference.
\end{document}
